\section{Nonlinearity and feature learning: proof of the additional theorem}
\label{sec:nonlinear-learning-proof}

This section proves \Cref{thm:nonlinear-learning} from the finite network
equations. All auxiliary letters and constants are local to their lemmas;
constants may depend on the fixed data, labels, activations and depth, but
not on width. Write $v_a=x_a/\sqrt d$. The extra hypotheses imply $d\ge2$.
We first detect input nonlinearity, then establish positive initial feature
accelerations, and finally control their real-time remainders. No limit of
the trained dynamics is used as an input.

\subsection{A finite witness of input nonlinearity}

\begin{lemma}[Initial velocity and four passive inputs]
\label{fl:input-lemma}
On each fixed finite query list, the initial prediction velocities converge
in probability to a deterministic continuous function on the sphere that is
not affine. Four deterministic sphere inputs and a linear combination of
their evaluations annihilate every affine function but not this limiting
velocity. They depend only on the fixed problem, not on realized weights or
later training. Every initial training-feature Gram is positive definite in
the width limit.
\end{lemma}

\begin{proof}
Gaussian covariance recursion gives a dot-product kernel
$F_\ell(v^\top u)$ at layer $\ell$, with $F_0(s)=s$.
Here is the needed expansion, including nonzero activation means.
If $G,G'$ are standard Gaussians of correlation $r$, the normalized Hermite
polynomials satisfy
$\E[H_j(G)H_k(G')]=\mathbf1_{j=k}r^k$.
Compare coefficients in
$\E[e^{zG-z^2/2}e^{wG'-w^2/2}]=e^{rzw}$ to obtain this identity.
They are complete in Gaussian $L^2$: a function orthogonal to every
polynomial has entire transform $\E[h(G)e^{zG}]$ with all derivatives
zero, hence zero Fourier transform and is zero. Fourier uniqueness here
follows, for example, by Gaussian convolution, Fourier inversion and the
approximate-identity limit.

For any positive variance $q$, linear growth gives the expansion
$\phi_\ell(\sqrt qG)=\sum_k\alpha_kH_k(G)$ in $L^2$ and therefore
\[
 \E[\phi_\ell(\sqrt qG)\phi_\ell(\sqrt qG')]
 =\sum_{k\ge0}\alpha_k^2r^k,
 \qquad\sum_k\alpha_k^2<\infty.
\]
Its positive support is unbounded: finite support would make the continuous
activation a polynomial, and its bounded derivative would make it affine.
The covariance identity passes from finite sums by Cauchy--Schwarz,
including $r=\pm1$. At layer $\ell$, compose this series with
$F_{\ell-1}(s)/F_{\ell-1}(1)$, using the positive variance
$q=F_{\ell-1}(1)$. All coefficients remain nonnegative and their sum is
the finite next Gaussian second moment. Tonelli's theorem first at $s=1$
then absolute convergence justify regrouping even when the inner constant
term is nonzero. A positive inner coefficient of degree $j\ge1$ and an
outer coefficient of degree $k$ give a positive coefficient of degree $jk$.
Induction yields
\begin{equation}
 F_\ell(s)=\sum_{k\ge0}c_{\ell,k}s^k,
 \qquad c_{\ell,k}\ge0,\quad
 \sum_kc_{\ell,k}<\infty,
 \label{fl:positive-series}
\end{equation}
with unbounded positive support and uniform absolute convergence on
$[-1,1]$. Since $[(v_a^\top v_b)^k]_{a,b}\to I_m$ as $k\to\infty$,
and every such matrix is the Gram of the tensor features $v_a^{\otimes k}$,
\eqref{fl:positive-series} also makes every initial training-feature Gram
positive definite.

For the stronger non-affinity conclusion let $P$ denote orthogonal
projection onto affine functions on the unit sphere, with normalized
surface measure $\sigma$:
\[
 (Ph)(v)=\int h(u)\,d\sigma(u)
       +d\,v^\top\int u h(u)\,d\sigma(u).
\]
Projection in both variables of $(v^\top u)^k$ gives the positive kernel
\[
 R_k(v,u)=(v^\top u)^k-a_k-b_kv^\top u,
 \quad a_k=\int(z^\top u)^k\,d\sigma(z),\quad
 b_k=d\int(z^\top u)^{k+1}\,d\sigma(z).
\]
Indeed the projected kernel is the Gram of the coordinatewise projected
tensor feature $(I-P)v^{\otimes k}$. Rotational invariance gives the
displayed formula and makes $a_k,b_k$ independent of unit $u$.
Since $d\ge2$, dominated convergence gives $a_k,b_k\to0$. Thus
$[R_k(v_a,v_b)]_{a,b}\to I_m$ and is positive definite for all large $k$.

The limiting initial velocity is
\[
 g(v)=\frac2m\sum_a y_aF_L(v^\top v_a).
\]
If $g$ were affine, applying $I-P$, evaluating at the training inputs and
pairing with $y$ would give
\[
 0=\frac2m\sum_k c_{L,k}
       \sum_{a,b}y_ay_bR_k(v_a,v_b)>0.
\]
Every summand is nonnegative, and a sufficiently large positive coefficient
makes one strictly positive since $y\ne0$. This is a contradiction.
The exchanges follow from uniform absolute convergence and boundedness of
$P$ in the uniform norm.

Some great circle witnesses non-affinity. Otherwise the antipodal average
of $g$ would be constant on every great circle, hence globally constant.
Subtract it and extend the remaining odd function homogeneously to
$\R^d$. Its restriction to every two-dimensional subspace would be linear,
making the extension additive and homogeneous on every pair of vectors,
and hence linear globally. On a circle witnessing non-affinity, interpolate
at three distinct points and take a fourth where the affine interpolant
fails. The resulting affine dependence, normalized and signed, gives
\begin{equation}
 \sum_{i=1}^4|c_i|=1,\qquad \sum_i c_i=0,\qquad
 \sum_i c_i u_i=0,\qquad \sum_i c_i g(u_i)>0.
 \label{fl:four-witness}
\end{equation}
Here the $u_i$ are unit vectors; their physical inputs are $\sqrt d\,u_i$.

Finally, finite-query initial covariance convergence is elementary.
Conditioned on the preceding layer, rows are independent Gaussian tuples
with the empirical preceding Gram as covariance. Linear growth and Gaussian
fourth moments bound the conditional variance of every averaged feature
product by $O(1/n)$ on bounded covariance sets. Its conditional mean is
continuous in the covariance, including singular matrices: couple through
the positive square root and use the Lipschitz activation and
Cauchy--Schwarz. Conditional Chebyshev and induction give the Gram limits.
The exact identity
$\dot f_n(0,x)=\frac2m\sum_a y_a h^{(L)}(0,x_a)^\top h^{(L)}(0,x)/n$
then proves the asserted velocity convergence.
\end{proof}

\subsection{Uniform real-time control}

\begin{lemma}[A width-independent real bootstrap]
\label{fl:bootstrap}
There are deterministic $M,T>0$ such that, with probability tending to one,
the following bounds hold for $0\le t\le T$:
$\|w(t)\|/\sqrt n\le Mt$, hidden matrix increments are at most $Mt^2$
in Frobenius norm (divide the first-layer norm by $\sqrt n$), and all
preactivation and feature increments have neuron RMS at most $Mt^2$,
uniformly on the input sphere. Moreover
\begin{equation}
 \sup_{\|x\|=\sqrt d}|f_n(t,x)-t\dot f_n(0,x)|\le Mt^2.
 \label{fl:first-remainder}
\end{equation}
\end{lemma}

\begin{proof}
For explicit local choices take $B\ge8$ bounding every
$|\phi_\ell(0)|$, $|\phi_\ell'|$ and $|\phi_\ell''|$ on the real line.
With probability tending to one,
$\|W^{(1)}(0)\|_{\rm op}/\sqrt n\le B$ and
$\|W^{(\ell)}(0)\|_{\rm op}\le B$ for $\ell\ge2$.
The first assertion follows from the fixed-$d$ law of large numbers for
$W^{(1)\top}W^{(1)}/n$. For the others, two $1/4$ sphere nets and Gaussian
tails bound each failure probability by
$2\exp(2n\log9-nB^2/8)$.
Put
\[
 A=B+1,\quad H=(2A^2)^L,\quad D=A^{2L},\quad
 E=2Y^2DH^2,\quad T=\min\{1,(2YH\sqrt D)^{-1}\}.
\]
Bootstrap the scaled first operator and other hidden operator norms below
$A$. Forward induction bounds all feature RMS by $H$, uniformly on the
sphere; backward induction gives
$\|\delta_a^{(\ell)}\|/\sqrt n\le D\|w\|/\sqrt n$.
Energy dissipation gives $\|r\|/\sqrt m\le Y$, hence
$\|w(t)\|/\sqrt n\le2YHt$.
The exact updates and sample Cauchy--Schwarz bound the scaled first
Frobenius speed by $2YD\|w\|/\sqrt n$ and later Frobenius speeds by
$2YDH\|w\|/\sqrt n$. Their increments are at most $Et^2$.
Since $ET^2\le1/2$, the bootstrap closes with strict margin.

Forward subtraction gives the common feature increment bound $Jt^2$, where
\[
 J=BEH\sum_{j=0}^{L-1}(BA)^j.
\]
The corresponding preactivation bound
follows from the same recursion before applying the Lipschitz activation.
Also $|f_n(t,x)|\le2YH^2t$ and
\[
 \frac{\|\dot w(t)-\dot w(0)\|}{\sqrt n}
 \le4YH^3t+2YJt^2.
\]
Integrate and write the output remainder as
$(w-t\dot w(0))^\top h^{(L)}(0)/n
 +w^\top(h^{(L)}(t)-h^{(L)}(0))/n$.
Its absolute value is at most
$[2YH^4+(8/3)YHJT]t^2$.
Choose $M$ to dominate these finitely many constants. All bounds use only
real derivatives, not a complex-time radius.
\end{proof}

\subsection{Positive initial accelerations}

\begin{lemma}[Finite initialization and layerwise activity]
\label{fl:initial-activity}
Under the additional theorem's hypotheses, each layer's initialized forward
fields, first backward derivatives (also before activation gates), and
preactivation and feature accelerations have deterministic joint empirical
row-law limits with convergent second moments. This includes the full
first-layer weight and acceleration rows. Every hidden block's squared
acceleration in inverse-mobility norm and every layer's squared feature
acceleration in neuron RMS converge to strictly positive constants.
The corresponding squared norms are uniformly integrable.
\end{lemma}

\begin{proof}
All notation below is local and all unmarked network fields are at zero.
Put $k=2/m$,
$\|q\|_n=\|q\|_F/\sqrt n$ for any array with $n$ rows, and define
\begin{align}
 s&=\sum_a y_a h_a^{(L)},&
 B_a^{(L)}&=s\odot\phi_L'(z_a^{(L)}),\nonumber\\
 B_a^{(\ell)}&=\phi_\ell'(z_a^{(\ell)})\odot
                  W_0^{(\ell+1)\top}B_a^{(\ell+1)},
 &V_\ell&=\ddot W^{(\ell)}(0),\nonumber\\
 R_a^{(\ell)}&=\ddot z_a^{(\ell)}(0),&
 A_a^{(\ell)}&=\ddot h_a^{(\ell)}(0).
 \label{fl:initial-fields}
\end{align}
Here $W_0^{(\ell)}=W^{(\ell)}(0)$, and
$\dot\delta_a^{(\ell)}(0)=kB_a^{(\ell)}$.
Each layer's joint empirical row law of its forward fields, $B,R,A$, and
the pre-gated fields $s$ or $W_0^{(\ell+1)\top}B_a^{(\ell+1)}$
converges in probability in quadratic Wasserstein distance to a deterministic
law. At layer one this includes the whole root row
$g_i=W_{0,i:}^{(1)}\in\R^d$ and the acceleration row $V_{1,i:}$.
In particular, all same-layer normalized inner products converge.

For the mobility norms
$\|V_1\|_{\rm mob}=\|V_1\|_n$ and
$\|V_\ell\|_{\rm mob}=\|V_\ell\|_F$ for $\ell\ge2$, there are
deterministic $e_\ell,c_\ell>0$ such that
\begin{equation}
 \|V_\ell\|_{\rm mob}^2\longrightarrow e_\ell,
 \qquad \sum_a\|A_a^{(\ell)}\|_n^2\longrightarrow c_\ell
 \quad\hbox{in probability}.
 \label{fl:positive-accelerations}
\end{equation}
We also establish uniform integrability of these squared norms and uniform
integrability in probability of their squared coordinate tails.
Since all hidden velocities initially vanish, differentiating the physical
updates and the forward recursion gives exactly
\begin{align}
 V_1&=k^2\sum_a y_aB_a^{(1)}v_a^\top,&
 V_\ell&=\frac{k^2}{n}\sum_a y_aB_a^{(\ell)}
                    h_a^{(\ell-1)\top}\quad(\ell\ge2),\nonumber\\
 R_a^{(1)}&=V_1v_a,&
 R_a^{(\ell)}&=V_\ell h_a^{(\ell-1)}
                         +W_0^{(\ell)}A_a^{(\ell-1)},\nonumber\\
 A_a^{(\ell)}&=\phi_\ell'(z_a^{(\ell)})\odot R_a^{(\ell)}.
 \label{fl:acceleration-recursion}
\end{align}
We prove the needed row-law limits by a finite forward--backward--forward
Gaussian conditioning argument, retaining both directions of each matrix.

For one hidden matrix $W$, suppose its recorded actions are $WH=Z$ and
$W^\top U=T$, with fixed column counts. On invertible query Grams set
$P_H=H(H^\top H)^{-1}H^\top$ and likewise $P_U$. Conditional on the
recorded transcript, its remaining law is
\begin{equation}
 W=Z(H^\top H)^{-1}H^\top
    +U(U^\top U)^{-1}T^\top(I-P_H)
    +(I-P_U)\widetilde W(I-P_H),
 \label{fl:two-sided-conditioning}
\end{equation}
where $\widetilde W$ has independent $N(0,1/n)$ entries. Indeed
$U^\top Z=T^\top H$, so the displayed mean satisfies both constraints.
Their homogeneous solution space consists exactly of
$(I-P_U)E(I-P_H)$, and that mean is Frobenius-orthogonal to the space.
Independence of orthogonal Gaussian projections proves the formula.
Before a reverse call omit its absent query block.

The identity also holds for interlaced adaptive calls to the independent
matrices. Condition first on the first-layer roots. Given the answers so
far, the next query used here is fixed and reveals a linear projection of
one remaining Gaussian matrix subspace. Its orthogonal complement stays
Gaussian and independent of all other residual matrices. Induction on calls
proves the stated conditional law for the complete global transcript.

In particular, for the first reverse call write
$Q_n=H^\top H/n$, $C_n=Z^\top U/n$, $D_n=U^\top U/n$. It gives
\begin{equation}
 T\overset d=H Q_n^{-1}C_n+(I-P_H)\Xi D_n^{1/2},
 \qquad
 \E[\|P_H\Xi D_n^{1/2}\|_n^2\mid\mathrm{past}]
   =\frac{\operatorname{rank}H}{n}\operatorname{tr}D_n,
 \label{fl:reverse-conditioning}
\end{equation}
where $\Xi$ has independent standard Gaussian entries, independent of the
past. After this answer, a new forward query $A$ obeys
\begin{align}
 F_n&=H^\top A/n,&J_n&=T^\top A/n,&E_n&=A^\top A/n,\nonumber\\
 \Sigma_n&=E_n-F_n^\top Q_n^{-1}F_n,\nonumber\\
 WA&\overset d=ZQ_n^{-1}F_n
  +UD_n^{-1}(J_n-C_n^\top Q_n^{-1}F_n)
  +(I-P_U)\Xi\Sigma_n^{1/2}.
 \label{fl:acceleration-conditioning}
\end{align}
Now $\Xi$ is fresh after the reverse answer. Since
$\Sigma_n=A^\top(I-P_H)A/n\succeq0$, it may be singular and is never
inverted. The removed projection's conditional squared RMS is
$\operatorname{rank}(U)\operatorname{tr}(\Sigma_n)/n$.

Here are the complete limit rules for these finitely many calls. Appending
independent standard Gaussian rows to an array with a deterministic empirical
$\mathcal W_2$ limit appends independent Gaussian variables to that limit.
For bounded continuous tests, conditional variances of empirical averages
are $O(1/n)$ and their conditional means converge by the old empirical law.
The appended second moment converges by the Gaussian law of large numbers.
Weak convergence and convergence of second moments give $\mathcal W_2$
convergence. Same-layer empirical products therefore converge, since they
are continuous of at most quadratic growth. Gram inverses converge whenever
their limits are positive definite. Positive square roots converge even at
singular covariance limits: bounded subsequences have limits whose squares
equal the limiting covariance, and its positive square root is unique.
Thus multiplying the fresh Gaussian rows by these roots preserves joint
convergence. The projection errors above vanish in RMS in probability,
since their covariance traces stay bounded in probability. Coupling rows
by their indices shows that these errors do not change a $\mathcal W_2$ limit.

The coordinate maps are linear combinations, Lipschitz activations, and
$(z,q)\mapsto\phi_\ell'(z)q$. Each is continuous with at most linear
growth in its complete input tuple and therefore preserves $\mathcal W_2$
convergence: on an $L^2$ coupling, continuity gives convergence in probability,
and uniform integrability of squared inputs and the linear envelope give
$L^2$ convergence. Coefficients converging to constants are handled first
by subtracting their linear combinations and using bounded input RMS.

Start with the iid full rows $g_i\sim N(0,I_d)$. Their empirical laws
converge with second moments. Make all initial forward calls in order.
Conditional fresh Gaussian rows and the preceding rules yield the forward
tuple limits $Z^{(\ell)}\sim N(0,Q^{(\ell-1)})$ and
$H_a^{(\ell)}=\phi_\ell(Z_a^{(\ell)})$.
Lemma~\ref{fl:input-lemma} gives $Q^{(\ell)}\succ0$ for $\ell\ge1$.

Write $\mathsf B,\mathsf R,\mathsf A$ for limiting row variables of
$B,R,A$, and $D^{(\ell)}_{ab}=\E[\mathsf B_a^{(\ell)}
\mathsf B_b^{(\ell)}]$. The top response is the continuous, at-most-linear
map $\mathsf B_a^{(L)}=(\sum_b y_bH_b^{(L)})\phi_L'(Z_a^{(L)})$.
Its Gram is positive definite. If a linear combination with coefficients
$c_a$ vanishes, full Gaussian support and continuity imply
\[
 \Bigl(\sum_b y_b\phi_L(z_b)\Bigr)
 \Bigl(\sum_a c_a\phi_L'(z_a)\Bigr)=0
 \quad\hbox{for every }z\in\R^m.
\]
The first factor is not identically zero, since $y^\top Q^{(L)}y>0$.
The second vanishes on a nonempty open set, hence everywhere by analyticity.
Varying one coordinate and using the nonconstant derivative gives $c_a=0$.

Proceed downward. In \eqref{fl:reverse-conditioning} take
$H=[h_a^{(\ell)}]$, $Z=[z_a^{(\ell+1)}]$ and
$U=[B_a^{(\ell+1)}]$. All upper responses are already known from the
forward pass and earlier reverse calls; this is the first reverse call to
the present matrix. The preceding limit rules give its limiting answer
\[
 \sum_bH_b^{(\ell)}[(Q^{(\ell)})^{-1}C]_{ba}+\Gamma_a,
 \quad C_{ba}=\E[Z_b^{(\ell+1)}\mathsf B_a^{(\ell+1)}],
 \quad\Gamma\sim N(0,D^{(\ell+1)}),
\]
where $\Gamma$ is independent of the previous lower-layer tuple, including
$g$ at layer one. Multiplication by $\phi_\ell'(Z_a^{(\ell)})$ yields
$\mathsf B_a^{(\ell)}$, and conditional variance gives, for $c\ne0$,
\[
 c^\top D^{(\ell)}c\ge\lambda_{\min}(D^{(\ell+1)})
       \sum_a c_a^2\E[\phi_\ell'(Z_a^{(\ell)})^2]>0.
\]
Every marginal preactivation is nondegenerate Gaussian; a nonzero analytic
derivative has a discrete zero set. This proves the strict inequality and
completes both the backward row-law and positive-Gram inductions.

Construct the accelerations forward using \eqref{fl:acceleration-recursion}.
At layer one the joint limit retains the exact row formulas
\[
 Z_a^{(1)}=g^\top v_a,\quad
 V_{1,\mathrm{row}}=k^2\sum_a y_a\mathsf B_a^{(1)}v_a,\quad
 \mathsf R_a^{(1)}=V_{1,\mathrm{row}}^\top v_a,\quad
 \mathsf A_a^{(1)}=\phi_1'(Z_a^{(1)})\mathsf R_a^{(1)}.
\]
For later layers the learned-link contribution is exactly
$k^2\sum_b y_bB_b^{(\ell)}Q_{n,ba}^{(\ell-1)}$ and has its row-law
limit. Apply \eqref{fl:acceleration-conditioning} to the remaining term
with $A=[A_a^{(\ell-1)}]$ and the already recorded forward and reverse
queries. Its coefficients are same-layer second moments already obtained;
its only inverse Grams are the positive $Q^{(\ell-1)},D^{(\ell)}$.
This proves joint convergence of the new preactivation and feature
accelerations. There are exactly $3(L-1)$ hidden-matrix calls; each query
is determined by earlier answers. Inverses are used only on events of
probability tending to one, which suffices for these limits.

Uniform integrability follows without inverses. The variable
$M_n=1+\|W_0^{(1)}\|_n+\sum_{\ell\ge2}\|W_0^{(\ell)}\|_{\rm op}$
has bounded moments of every fixed order, uniformly in width. For the first
term use Jensen and Gaussian moments; for the others integrate the two-net
Gaussian tail used in Lemma~\ref{fl:bootstrap}. Linear activation growth,
bounded derivative gates and \eqref{fl:acceleration-recursion} bound every
array RMS and block mobility norm by a fixed polynomial in $M_n$.
Thus their squared norms have bounded moments above order one. This gives
uniform integrability of the energies, while the proved $\mathcal W_2$
convergence gives uniform integrability in probability of squared coordinate
tails. In particular all limits of squared norms are also $L^1$ limits.

Finally the exact transpose relation in \eqref{fl:initial-fields} gives
the partial-layer adjoint identity
\begin{equation}
 \sum_a\frac{y_a}{n}B_a^{(\ell)\top}R_a^{(\ell)}
       =\frac1{k^2}\sum_{j=1}^{\ell}\|V_j\|_{\rm mob}^2.
 \label{fl:partial-adjoint}
\end{equation}
Indeed the learned-link term is $\|V_\ell\|_{\rm mob}^2/k^2$ by
the formula for $V_\ell$, and the propagated term equals the left side
one layer lower; at the first layer only the learned-link term remains.
The squared mobility limits are
$e_\ell=k^4y^\top(Q^{(\ell-1)}\circ D^{(\ell)})y$.
For $Q=\sum_jq_jq_j^\top\succeq0$ and $D\succ0$,
$Q\circ D=\sum_j\operatorname{diag}(q_j)D\operatorname{diag}(q_j)
\succeq\lambda_{\min}(D)\operatorname{diag}(Q)$.
Every relevant $Q$ has positive diagonal, including $Q^{(0)}$, so
$e_\ell>0$. Passing to the established joint-law limits in
\eqref{fl:partial-adjoint} shows that some $\mathsf R_a^{(\ell)}$ is
nonzero in $L^2$. Since $\phi_\ell'(Z_a^{(\ell)})\ne0$ almost surely,
$\mathsf A_a^{(\ell)}=\phi_\ell'(Z_a^{(\ell)})\mathsf R_a^{(\ell)}$
cannot vanish for all $a$. Thus
$c_\ell=\sum_a\E|\mathsf A_a^{(\ell)}|^2>0$, as claimed.
\end{proof}


\subsection{From accelerations to actual feature learning}

\begin{lemma}[Strong early-time expansion]
\label{fl:strong-time}
Every hidden layer has the activity asserted in
\eqref{fl:layer-motion}. Moreover, on the training inputs,
\begin{equation}
 y^\top(f_n(t)-f_{\mathrm{NTK},n}(t))
 =\frac m3 t^3\left[
 \frac{\|\ddot W^{(1)}(0)\|_F^2}{n}
 +\sum_{\ell=2}^L\|\ddot W^{(\ell)}(0)\|_F^2\right]+o_*(t^3).
 \label{fl:cubic-identity}
\end{equation}
Here $o_*(t^k)$ means that for every $\varepsilon>0$ there is a
deterministic $T_\varepsilon>0$, independent of width, such that the
probability of a remainder exceeding $\varepsilon t^k$ anywhere on
$0<t\le T_\varepsilon$ tends to zero. The bracket converges in probability
to a strictly positive deterministic constant.
\end{lemma}

\begin{proof}
We give the remainder argument because differentiating at zero at each
width would not, by itself, prove a gap at a width-independent positive time.
Use the accelerations
$V^{(\ell)}=\ddot W^{(\ell)}(0)$,
$R_a^{(\ell)}=\ddot z_a^{(\ell)}(0)$,
$A_a^{(\ell)}=\ddot h_a^{(\ell)}(0)$, and write
$\|u\|_n=\|u\|_2/\sqrt n$.
Let $P_a^{(L)}=\sum_b y_bh_b^{(L)}(0)$ and
$P_a^{(\ell)}=W^{(\ell+1)}(0)^\top B_a^{(\ell+1)}$ below the top.
Recursively define
$B_a^{(\ell)}=\phi_\ell'(z_a^{(\ell)}(0))\od P_a^{(\ell)}$.
Thus $(2/m)B_a^{(\ell)}=\dot\delta_a^{(\ell)}(0)$, as in
\Cref{fl:initial-activity}.
For a tuple of hidden matrices use the local norm
\[
 \|U\|_{\rm mob}^2=\|U^{(1)}\|_F^2/n+
                       \sum_{\ell=2}^L\|U^{(\ell)}\|_F^2.
\]
The initialization lemma supplies deterministic joint limits with second
moments, including the pre-gated $P$ and the preactivation accelerations $R$.
Consequently, for every $\eta>0$ there is a finite $D$ such that
\begin{equation}
 \Pr\left\{\max_u\frac1n\sum_i |u_i|^2
                     \mathbf1_{|u_i|>D}>\eta\right\}\longrightarrow0,
 \label{fl:tail-control}
\end{equation}
where $u$ ranges over these finitely many initial vectors and their fixed
linear combinations. To verify this consequence, choose a continuity radius
$D$ at which each deterministic limiting tail second moment is below
$\eta/2$, then use convergence with second moments.

The only nonlinear remainder estimate needed is the following. For bounded
Lipschitz $g$, a fixed initial vector $u$, and any $z,z_0$,
\begin{equation}
 \|[g(z)-g(z_0)]\od u\|_n^2
 \le \operatorname{Lip}(g)^2D^2\|z-z_0\|_n^2
       +4\|g\|_\infty^2\|u\mathbf1_{|u|>D}\|_n^2.
 \label{fl:multiplier}
\end{equation}
Split the coordinates at $|u_i|=D$ and use, respectively, the Lipschitz
and boundedness estimates. By \Cref{fl:bootstrap}, the first term is
bounded by a deterministic constant times $D^2t^4$ on an event whose
probability tends to one. Choose $D$ using \eqref{fl:tail-control}, then
choose a sufficiently small deterministic time. This proves an $o_*(1)$
bound, also uniformly when $z-z_0$ is replaced by $\theta(z-z_0)$,
$0\le\theta\le1$. This order of choices is what gives the stated
probability quantifier, rather than only a confidence-dependent time.

The readout equation and the bootstrap give
\[
 \left\|w(t)/t-\dot w(0)\right\|_n\le C_0t,
 \qquad \dot w(0)=\frac2m\sum_b y_bh_b^{(L)}(0),
\]
where $C_0$ is local and may depend on the entire fixed problem.
At the top, subtract the backward recursion from its proposed leading term
$(2t/m)B_a^{(L)}$: the difference divided by $t$ is the bounded gate
times $w(t)/t-\dot w(0)$, plus a gate difference times $\dot w(0)$.
The latter is controlled by \eqref{fl:multiplier}. At a lower layer,
writing $D_a^{(\ell)}(t)=\operatorname{diag}
(\phi_\ell'(z_a^{(\ell)}(t)))$, the exact difference divided by $t$ is
\begin{align*}
 &D_a^{(\ell)}(t)W^{(\ell+1)}(t)^\top
       [\delta_a^{(\ell+1)}(t)/t-(2/m)B_a^{(\ell+1)}]\\
 &\quad+\frac2mD_a^{(\ell)}(t)
       [W^{(\ell+1)}(t)-W^{(\ell+1)}(0)]^\top B_a^{(\ell+1)}
       +\frac2m[D_a^{(\ell)}(t)-D_a^{(\ell)}(0)]P_a^{(\ell)}.
\end{align*}
The first term propagates the next layer's $o_*(1)$ error through bounded
operators, the second is $O(t^2)$ in RMS, and the third is $o_*(1)$ by
\eqref{fl:multiplier}. Downward induction and then the weight equations give
\begin{equation}
 \delta_a^{(\ell)}(t)=\frac{2t}{m}B_a^{(\ell)}+o_{*,\|\cdot\|_n}(t),
 \qquad
 W(t)-W(0)=\frac{t^2}{2}V+o_{*,\|\cdot\|_{\rm mob}}(t^2).
 \label{fl:weight-expansion}
\end{equation}
For the second assertion, use $r=-y+O(t)$ and $h-h(0)=O_{\|\cdot\|_n}(t^2)$.
The exact identities
$\|uv^\top\|_F/\sqrt n=\|u\|_n\|v\|_2$ and
$\|uh^\top\|_F/n=\|u\|_n\|h\|_n$ bound each difference of outer
products. They give $\dot W-tV=o_{*,\|\cdot\|_{\rm mob}}(t)$;
integrating gives \eqref{fl:weight-expansion} with its uniform quantifier.

For the forward pass, the first preactivation expansion follows directly
from the first matrix expansion. If
$\Delta z=t^2R/2+o_{*,\|\cdot\|_n}(t^2)$, the exact integral identity
\[
 \phi(z_0+\Delta z)-\phi(z_0)
    =\int_0^1\phi'(z_0+\theta\Delta z)\od\Delta z\,d\theta
\]
has remainder, after subtracting $t^2\phi'(z_0)\od R/2$, equal to
\[
 \int_0^1\phi'(z_0+\theta\Delta z)\od(\Delta z-t^2R/2)\,d\theta
 +\frac{t^2}{2}\int_0^1[\phi'(z_0+\theta\Delta z)-\phi'(z_0)]\od R\,d\theta.
\]
Both are $o_{*,\|\cdot\|_n}(t^2)$ by bounded gates and
\eqref{fl:multiplier}. In the next layer use
$\Delta z=\Delta W h(0)+W(0)\Delta h+\Delta W\Delta h$;
the last term is $O_{\|\cdot\|_n}(t^4)$ and the others have the
initialized acceleration coefficient. Thus
\begin{equation}
 h_a^{(\ell)}(t)-h_a^{(\ell)}(0)
    =\frac{t^2}{2}A_a^{(\ell)}+o_{*,\|\cdot\|_n}(t^2)
 \label{fl:feature-expansion}
\end{equation}
at every layer. The positive limiting squared acceleration norms from
\Cref{fl:initial-activity} prove \eqref{fl:layer-motion} on one common
small interval. In particular, a nonzero weight increment has not been
mistaken for a nonzero feature increment. Nor have we assumed that a
nonlinear activation defines a Fr\'echet-smooth map on $L^2$.

For the output gap, the exact mobility-weighted tangent kernel is
\begin{align*}
 K_{ab}(t)={}&\langle h_a^{(L)}(t),h_b^{(L)}(t)\rangle_n
 +(v_a^\top v_b)\langle\delta_a^{(1)}(t),\delta_b^{(1)}(t)\rangle_n\\
 &+\sum_{\ell=2}^L
  \langle\delta_a^{(\ell)}(t),\delta_b^{(\ell)}(t)\rangle_n
  \langle h_a^{(\ell-1)}(t),h_b^{(\ell-1)}(t)\rangle_n,
 \qquad \langle u,v\rangle_n=u^\top v/n.
\end{align*}
It satisfies $\dot f_n=(2/m)K(y-f_n)$ and $K(0)=Q_n^{(L)}(0)$.
Equations \eqref{fl:weight-expansion}--\eqref{fl:feature-expansion} yield
$K(t)=K(0)+t^2K_2+o_*(t^2)$. Its readout contribution to $K_2$ is
$\tfrac12(\langle A_a^{(L)},h_b^{(L)}(0)\rangle_n+
\langle h_a^{(L)}(0),A_b^{(L)}\rangle_n)$; the hidden contributions
replace each $\delta$ by $(2/m)B$ and each feature by its initial value.
All coefficient norms are bounded on events whose probability tends to one.

Write $E_n=\|V\|_{\rm mob}^2$. The partial-layer adjoint identity in
the initialization lemma, at the top, gives
\[
 \left\langle\sum_a y_ah_a^{(L)}(0),\sum_b y_b A_b^{(L)}\right\rangle_n
   =\frac{m^2}{4}E_n.
\]
This is exactly the readout part of $y^\top K_2y$. Expanding the squared
outer-product formulas for $V$ shows that the hidden part is also
$m^2E_n/4$. Therefore
\begin{equation}
 y^\top K_2y=\frac{m^2}{2}E_n.
 \label{fl:positive-kernel}
\end{equation}
Let $e=f_n-f_{\mathrm{NTK},n}$ on the training set. Subtract the two
prediction equations and integrate the constant-kernel part exactly:
\[
 e(t)=\frac2m\int_0^t e^{-(2/m)K(0)(t-s)}
                  [K(s)-K(0)][y-f_n(s)]\,ds.
\]
The exponential is contractive since $K(0)\succeq0$, and its difference
from the identity is $O(t-s)$ on the bounded coefficient event.
The kernel remainder contributes $o_*(t^3)$; the factors $f_n(s)=O(s)$
and this exponential difference contribute $O(t^4)$. Hence
$e(t)=(2t^3/3m)K_2y+o_*(t^3)$, proving
\eqref{fl:cubic-identity}. Since $E_n$ has a positive deterministic limit,
this also proves a positive label-projected gap on one deterministic
small interval, with probability tending to one.
\end{proof}

\subsection{The changing features are genuinely nonlinear}

\begin{lemma}[Nonlinear first-layer increment]
\label{fl:nonlinear-feature}
Let $V=\operatorname{span}\{x_1,\ldots,x_m\}$ and let $\sigma_V$ be
normalized surface measure on its unit sphere. For all sufficiently small
fixed $t>0$, with probability tending to one,
\begin{equation}
 \frac1n\sum_{i=1}^n\inf_{a\in V,\,b\in\R}
 \int_{S(V)}\bigl|h_i^{(1)}(t,\sqrt d\,v)-h_i^{(1)}(0,\sqrt d\,v)
                         -a^\top v-b\bigr|^2\,d\sigma_V(v)
 \ge c_*t^4.
 \label{fl:nonlinear-increment}
\end{equation}
Here $c_*>0$ depends on the fixed problem, not on width.
\end{lemma}

\begin{proof}
The nonparallel-input assumption implies $\dim V\ge2$. For nonzero
$g,r\in V$, the function
$v\mapsto\phi_1'(g^\top v)(r^\top v)$ is not affine on $S(V)$.
Indeed, if it were $a^\top v+b$, evaluating at the equator $r^\perp$
and its antipodes gives $b=0$ and $a=\lambda r$. Away from that equator,
$\phi_1'(g^\top v)=\lambda$. Continuity extends this to the whole
sphere, so $\phi_1'$ is constant on $[-\|g\|,\|g\|]$. Real analyticity
would then make $\phi_1$ affine on $\R$, a contradiction.

Project the initialized first-layer row onto $V$ and call it $g_i$;
call its acceleration row $r_i$. The latter already lies in $V$ by its
explicit update formula. The initialization lemma gives a deterministic
joint empirical limit with second moments: $g$ is a nondegenerate Gaussian
on $V$, and $\E\|r\|^2>0$. Thus $g\ne0$ almost surely and $r\ne0$
with positive probability. Let $P$ now denote the orthogonal projection
onto affine functions in $L^2(S(V),\sigma_V)$. The continuous function
\[
 (g,r)\longmapsto
 \|(I-P)[\phi_1'(g^\top v)(r^\top v)]\|_{L^2(\sigma_V)}^2
\]
is bounded by $\|\phi_1'\|_\infty^2\|r\|^2$ and is strictly positive
when both arguments are nonzero. Joint convergence with second moments
therefore makes its empirical average converge to a deterministic
positive constant.

By \eqref{fl:weight-expansion}, the first-layer row increment is
$t^2r_i/2+o_*(t^2)$ in row RMS. A bounded activation derivative transfers
this error to $L^2$ of neuron index times $\sigma_V$. For the remaining
fixed-direction Taylor expansion, truncate at $\|r_i\|\le D$.
There the remainder divided by $t^2$ is bounded by
$\|\phi_1''\|_\infty t^2D^2/8$. On the complement it is bounded
by $\|\phi_1'\|_\infty\|r_i\|$, whose RMS tail tends to zero by
the initialized joint second-moment limit. Choosing $D$ and then time
as in \eqref{fl:tail-control} gives an $o_*(t^2)$ remainder in the
product space. Projection $I-P$ is contractive. The positive limiting
projected acceleration norm consequently proves
\eqref{fl:nonlinear-increment}. This argument needs no full input rank
and permits unbounded activation values.
\end{proof}

\subsection{Transfer to all three compressed predictors}

\begin{proof}[Proof of \Cref{thm:nonlinear-learning}]
First consider the dense reference. Use the four inputs and normalized
coefficients of \eqref{fl:four-witness}; denote their positive limiting
velocity combination by $a_0$. Initial velocity convergence and
\eqref{fl:first-remainder} imply, on a sufficiently small deterministic
interval and with probability tending to one,
\[
 \sum_{i=1}^4c_i f_n(t,\sqrt d\,u_i)\ge a_0t/2.
\]
Every affine predictor is annihilated by this combination. Since
$\sum_i|c_i|=1$, its distance from $f_n(t,\cdot)$ on any query set
containing these four points is at least $a_0t/2$.
Similarly, \eqref{fl:cubic-identity} and the positive energy limit give
$y^\top(f_n-f_{\mathrm{NTK},n})\ge c_0t^3$ on the training set,
so its largest coordinate gap is at least $c_0t^3/\|y\|_1$.
The layer and nonlinear-increment conclusions are
\Cref{fl:strong-time,fl:nonlinear-feature}.

Now fix any such positive time before taking $n\to\infty$.
The compression theorem gives
$\|f_{\rm model}-f_n\|_{\mathcal X}\to0$ in probability.
For completeness, Harmonic and Taylor have the displayed eventual
$Y/n$ absolute estimate at every confidence, which implies this convergence;
Legendre has the vanishing absolute estimate used in its headline assembly.
The triangle inequality transfers both dense gaps to the compressed
predictor, with half their constants, on events whose probabilities tend
to one. Take a common smaller $c_*,t_*$ for the finite list of conclusions.

Legendre and Harmonic already promise accuracy on the whole sphere. Taylor
must include the four witnesses as passive inputs before initialization;
their selection uses the deterministic covariance recursion and labels,
not realized weights or later training, and supplies no passive labels.
Its proved bound then uses $p+4$ instead of $p$. Since $m+p\ge2$,
this increases the displayed quadratic panel factor by at most $9$ and
the linear data factor by at most $3$, without changing any exponent.
An arbitrary unaugmented panel need not witness input non-affinity.
No assertion is made about correspondence between individual compressed
coordinates and dense features, nor about a gap at the fitted endpoint.
\end{proof}
